\title{xLSTM: Extended Long Short-Term Memory}

\begin{abstract}
During the 1990s, the foundational concepts of the Long Short-Term Memory (LSTM) model—namely, the constant error carousel and gating mechanisms—were established. Over the ensuing decades, LSTMs have remained relevant and have played a pivotal role in significant advances across deep learning, particularly serving as the original Large Language Models (LLMs). Nevertheless, the introduction of Transformers, which utilize self-attention mechanisms allowing for efficient parallelization, heralded a transformative period in the field, ultimately outperforming LSTMs at large scales. Motivated by this, we investigate the limits of language modeling performance by scaling LSTMs to the order of billions of parameters, incorporating insights from recent LLM developments while directly addressing previously identified LSTM shortcomings. Our approach introduces exponential gating in conjunction with advanced normalization and stabilization strategies. Further, we alter the internal LSTM memory configuration, resulting in two variants: (i) sLSTM, characterized by scalar memory elements, scalar-based updates, and a novel form of memory mixing, and (ii) mLSTM, a fully parallel version equipped with matrix-based memory and an update procedure grounded in covariance computations. By integrating these LSTM modifications within residual block frameworks, we form xLSTM blocks, which are subsequently aggregated into xLSTM architectures via residual stacking. These improvements in gating and memory formulation significantly enhance the capabilities of xLSTM, enabling it to achieve results that rival those of advanced Transformers and State Space Models in both scalability and performance.\\
Code available at: \url{https://github.com/NX-AI/xlstm}
\end{abstract}

\section{Introduction}
\label{sec:intro}

The Long Short-Term Memory (LSTM) ideas~\citep{Hochreiter:91,Hochreiter:97nips,Hochreiter:97}, i.e., the constant error carousel and gating, were introduced to overcome the vanishing gradient problem of recurrent neural networks~\citep{Hochreiter:91,Hochreiter:00book}:

\begin{align}
\label{eq:lstm_idea}
  c_t \ = \ f_t \ c_{t-1} \ + \ 
          i_t \  z_t \ , \quad  
          h_t \ = \ o_t \ \psi({c_t}) \ . 
\end{align}

The constant error carousel functions by incrementally adjusting the cell state $c_{t-1}$ (highlighted in green) through contributions from cell inputs $z_t$, with the adjustments being regulated by sigmoid gates (depicted in blue). The input gate $i_t$ together with the forget gate $f_t$ determine how the cell state is updated, while the output gate $o_t$ governs the release of information from the memory cell, corresponding to the hidden state $h_t$. Prior to generating the hidden state, the cell state undergoes normalization or compression through the $\psi$ operation, after which the output gate applies its influence to produce the final hidden state.

LSTMs have achieved significant success across a range of fields \citep{Hochreiter:01,Hochreiter:07,Schmidhuber:15} and maintained dominance in text generation until the introduction of Transformers in 2017~\citep{Vaswani:17}. Their capacity for handling sequences has been validated through diverse applications, including text generation~\citep{Graves:13,Karpathy:15blog}, handwriting synthesis~\citep{Graves:13}, sequence-to-sequence translation~\citep{Sutskever:14nips}, evaluation of computer code~\citep{Zaremba:14arxiv}, image captioning~\citep{Karpathy:15,Hossain:19}, code generation~\citep{Karpathy:15blog}, modeling rainfall-runoff processes~\citep{Kratzert:18,Kratzert:19}, and constructing hydrological models for flood warnings~\citep{Nearing:24}. Within reinforcement learning, LSTMs represent the leading architecture for sequence modeling, as demonstrated by systems such as AlphaStar for StarCraft II~\citep{Vinyals:17short}, OpenAI Five for Dota 2~\citep{OpenAI:19}, and the magnetic controller employed in nuclear fusion~\citep{Degrave:22short}. A notable strength of LSTMs lies in their ability to develop abstract representations, effectively extracting semantic features for storage in their internal memory cells~\citep{Karpathy:15blog}. This proficiency is highlighted by findings such as the discovery of number and syntax neurons~\citep{Lakretz:19}, linguistic neurons~\citep{Bau:19}, and sentiment neurons~\citep{Radford:17arxiv}. LSTMs continue to find relevance in a variety of prominent applications~\citep{Degrave:22short,Nearing:24}, underscoring their enduring utility.

Although LSTMs have achieved remarkable results, they are hindered by three primary drawbacks: (i) Inability to adjust prior storage choices. This is illustrated by the {\it Nearest Neighbor Search} task (refer to Appendix~\ref{sec:appNearestNeighborSearch}), where a given reference vector necessitates sequential inspection of a sequence to locate and return the value attached to the most similar vector by sequence end. The left panel in Figure~\ref{fig:lstmProblems} presents the mean squared error for this task. LSTM has difficulty modifying a previously stored value if a better match appears later in the sequence; in contrast, our proposed xLSTM addresses this challenge by employing exponential gating. (ii) Restricted memory capacity, as information must be condensed into scalar-valued cell states. We highlight this limitation through {\it Rare Token Prediction}. The right panel of Figure~\ref{fig:lstmProblems} displays token prediction perplexity on Wikitext-103~\citep{Merity:17} segmented by token frequency. LSTM underperforms for rare tokens owing to its storage constraints. Our xLSTM overcomes this shortcoming by introducing a matrix-based memory. (iii) Inability to parallelize, stemming from hidden-to-hidden state dependencies between consecutive time steps, which demand strict sequential computation.

These constraints of LSTM have contributed to the rise of the Transformer architecture~\citep{Vaswani:17} for language modeling tasks. Therefore, an important question arises: what levels of performance can be obtained in language modeling if these deficiencies are mitigated and LSTMs are scaled to the sizes of current Large Language Models?

\section{Extended Long Short-Term Memory}
\label{sec:xLSTM}

To address certain shortcomings of the LSTM, the Extended Long Short-Term Memory (xLSTM) framework proposes two essential changes to the formulation in Equation~\eqref{eq:lstm_idea}. These changes consist of exponential gating and innovative memory organization. As a result, two new models are introduced within the LSTM family: (i) sLSTM (see Section~\ref{sec:sLSTM}), which features scalar-valued memory, scalar updates, and incorporates memory mixing, and (ii) mLSTM (see Section~\ref{sec:mLSTM}), utilizing matrix-based memory with an update mechanism rooted in the covariance (outer product) operation that facilitates complete parallelization. Exponential gating is a distinguishing trait of both the sLSTM and mLSTM. In order to achieve full parallelization, mLSTM omits memory mixing—that is, there are no recurrent hidden-to-hidden connections. Both variants, sLSTM and mLSTM, may be expanded to encompass multiple memory cells; the sLSTM applies memory mixing between these cells. Moreover, sLSTM is extendable to multiple heads, within which memory mixing is limited to the cells of each head, without inter-head mixing. Introducing multiple heads alongside exponential gating in sLSTM yields a novel method for memory mixing. In the case of mLSTM, having multiple heads is operationally identical to having multiple cells.

These variants can be embedded into residual blocks, thereby creating xLSTM blocks (see Section~\ref{sec:xLSTMblock}). By stacking such blocks residually, one constructs xLSTM architectures (refer to Section~\ref{sec:xLSTMarchitecure}). The overall composition of the xLSTM architecture and its elements is depicted in Figure~\ref{fig:xlstm_sketch}.

\subsection{Review of the Long Short-Term Memory}
\label{sec:orgLSTM}

The foundational concept behind LSTM~\citep{Hochreiter:91,Hochreiter:97nips,Hochreiter:97} established the scalar memory cell as a pivotal element for computation and data retention, effectively mitigating the issue of vanishing gradients~\citep{Hochreiter:91,Hochreiter:00book} via the constant error carousel mechanism, which governs updates to the cell state. The memory cell features three gating components: input, output, and forget gates, with the forget gate later proposed by \citet{Gers:00}. The memory cell’s update equations at timestep $t$ are as follows:

\begin{align}
c_t \ &= \  f_t \ c_{t-1} \ + \ i_t \ z_t &  & &\text{cell state} \\
h_t  \ &= \ o_t \ \tilde{h}_t \ , & \tilde{h}_t \ &= \ \psi \left( c_t\right)
  &\text{hidden state} \\
z_t \ &= \ \varphi \left( \tilde{z}_t \right) \ , 
  &\tilde{z}_t \ &=  \ \mathbf{w}^\top_{z} \ \mathbf{x}_t \ + \
  r_{z}  h_{t-1} \ + \  b_{z} \ \
  &\text{cell input} \\
i_t \ &= \ \sigma \left( \tilde{i}_t  \right) \ , 
  &\tilde{i}_t \ &= \ \mathbf{w}^\top_{i} \ \mathbf{x}_t \ + \
  r_{i}  \ h_{t-1} \ + \  b_{i} \ \
  &\text{input gate} \\
f_t \ &= \ \sigma \left(  \tilde{f}_t \right) \ , 
  &\tilde{f}_t \ &= \ \mathbf{w}^\top_{f} \ \mathbf{x}_t  \ + \
  r_{f}  \ h_{t-1} \ + \  b_{f} \ \
  &\text{forget gate} \\
o_t \ &= \ \sigma \left( \tilde{o}_t \right) \ , 
  &\tilde{o}_t  \ &= \ \mathbf{w}^\top_{o} \ \mathbf{x}_t \ + \
  r_{o}  \ h_{t-1} \ + \  b_{o} \ \
  &\text{output gate} 
\end{align}

The weight vectors $\mathbf{w}_{z}$, $\mathbf{w}_{i}$,
$\mathbf{w}_{f}$, and $\mathbf{w}_{o}$ denote the input weights connecting the inputs $\mathbf{x}_t$ to the cell input, input gate, forget gate, and output gate, respectively. Similarly, the weights $r_{z}$, $r_{i}$, $r_{f}$, and $r_{o}$ represent the recurrent connections from the previous hidden state $h_{t-1}$ to the cell input and each gate. The bias parameters $b_{z}$, $b_{i}$, $b_{f}$, and $b_{o}$ correspond to their respective bias terms. The cell input and hidden state activation functions are given by $\varphi$ and $\psi$, with $\tanh$ typically used in both roles. Specifically, $\psi$ serves to scale the cell state to bounded values, preventing it from diverging. All gates utilize the sigmoid activation function, defined as $\sigma(x) = 1/(1+\exp(-x))$. Subsequent developments grouped several scalar memory cells $\mathbf{c}_t \in \mathbb{R}^{d}$ into a single vector representation, enabling recurrent weight matrices $\mathbf{R} \in \mathbb{R}^{d \times d}$ that facilitate interactions among memory cell outputs~\citep{Greff:15}; further elaboration is available in Appendix~\ref{sec:apporgLSTM}. Empirical ablation analyses indicate that each element of the memory cell is essential for optimal performance~\citep{Greff:15}.

\subsection{sLSTM}
\label{sec:sLSTM}

To enable LSTMs to adjust their storage choices, we propose the addition of exponential gates (highlighted in red) alongside normalization and stabilization mechanisms. Specifically, we allow the input and forget gates to utilize exponential activation functions. For normalization purposes, we present a normalizer state, which accumulates the result of multiplying the input gate by every subsequent forget gate.

The scalar sLSTM forward pass is:

\begin{align}
\label{eq:slstmforward}
c_t \ &= \  f_t \ c_{t-1} \ + \ i_t \ z_t & & 
 &\text{cell state} \\
n_t \ &= \  f_t \ n_{t-1} \ + \ i_t  & & 
 &\text{normalizer state} \\
h_t \ &= \ o_t \ \tilde{h}_t \ ,
  &\tilde{h}_t \ &= \ c_t / n_t
&\text{hidden state} \\
z_t \ &= \ \varphi \left( \tilde{z}_t \right) \ , 
  &\tilde{z}_t \ &=  \ \mathbf{w}^\top_{z} \ \mathbf{x}_t \ + \
  r_{z}  h_{t-1} \ + \  b_{z}
  &\text{cell input} \\
i_t \ &= \ \!\exp \left( \tilde{i}_t  \right) \ , 
  &\tilde{i}_t \ &= \ \mathbf{w}^\top_{i} \ \mathbf{x}_t \ + \
  r_{i}  \ h_{t-1} \ + \  b_{i}
  &\text{input gate} \\
f_t \ &= \ \sigma \left(  \tilde{f}_t \right) \ \text{OR} \
\!\exp \left(  \tilde{f}_t \right) \ ,
  &\tilde{f}_t \ &= \ \mathbf{w}^\top_{f} \ \mathbf{x}_t  \ + \
  r_{f}  \ h_{t-1} \ + \  b_{f}
  &\text{forget gate} \\
o_t \ &= \ \sigma \left( \tilde{o}_t \right) \ , 
  &\tilde{o}_t  \ &= \ \mathbf{w}^\top_{o} \ \mathbf{x}_t \ + \
  r_{o}  \ h_{t-1} \ + \  b_{o}
  &\text{output gate} 
\end{align}

We transfer the original LSTM gating techniques, i.e., input- and/or hidden-dependent gating plus bias term, to the new architectures. Exponential activation functions can lead to large values that cause overflows. Therefore, we stabilize gates with an additional state \mbox{$m_t$~\citep{Milakov:18arxiv}}:

\begin{align}
\label{eq:slstmstabil}
m_t \ &= \ \max \left( \log ( f_t ) + m_{t-1} , \log ( i_t ) \right) &\text{stabilizer state} \\
i'_t \ &= \ \!\exp \left( \log \left ( i_t \right) - m_t \right) \ = \ \!\exp \left( \tilde{i}_t - m_t \right) \ \, 
  &\text{stabil. input gate} \\
f'_t \ &= \ \!\exp \left( \log \left( f_t \right) + m_{t-1} - m_t \right) \ \, 
  &\text{stabil. forget gate}
\end{align}

In Appendix~\ref{sec:appsLSTM}, we demonstrate that substituting $f_t$ with $f'_t$ and $i_t$ with $i'_t$ during the forward computation leaves both the overall network output and the gradients of the loss relative to the parameters unaffected.

\paragraph{Emergent Memory Mixing.} Similar to traditional LSTM architectures, sLSTM is also capable of accommodating several memory cells (refer to Appendix~\ref{sec:appsLSTM} for details). Having multiple memory cells allows for interactions among memories through the recurrent links $\mathbf{R}_{\mathbf{z}}$, $\mathbf{R}_{\mathbf{i}}$, $\mathbf{R}_{\mathbf{f}}$, $\mathbf{R}_{\mathbf{o}}$; these originate from the hidden state $\mathbf{h}$ and connect to the memory cell input $\mathbf{z}$ and the gating units $\mathbf{i}$, $\mathbf{f}$, $\mathbf{o}$, respectively. A distinctive feature in memory mixing for sLSTM is the influence of exponential gating mechanisms. The revised sLSTM model can incorporate several heads, wherein memory mixing occurs independently within each head, without inter-head mixing. Combining heads and exponential gating in sLSTM opens up novel possibilities for organizing and mixing memory representations.

\subsection{mLSTM}
\label{sec:mLSTM}

To improve the memory capacity of LSTMs, we replace the traditional scalar memory cell $c \in \mathbb{R}$ with a matrix memory cell $\mathbf{C} \in \mathbb{R}^{d \times d}$. This modification enables memory retrieval through matrix multiplication. At a given time step $t$, the network stores a pair of vectors: a key $\mathbf{k}_t \in \mathbb{R}^d$ and a value $\mathbf{v}_t \in \mathbb{R}^d$ (drawing terminology from the Transformer architecture). Subsequently, at time $t + \tau$, the value $\mathbf{v}_t$ is accessed using a query vector $\mathbf{q}_{t+\tau} \in \mathbb{R}^d$. This framework is aligned with the paradigm of Bidirectional Associative Memories (BAMs) \citep{Kohonen:72,Anderson:72,Nakano:72,Anderson:77}.

The covariance update rule~\citep{Sejnowski:77,Dayan:91} for storing a key-value pair is

\begin{align}
\mathbf{C}_t \ &= \ \mathbf{C}_{t-1} \ + \ \mathbf{v}_{t} \ \mathbf{k}_t^\top \ .
\end{align}

We posit the presence of a layer normalization applied prior to the projection of inputs onto keys and values, ensuring their means are zero. The procedure for updating the covariance has been demonstrated to be optimal~\citep{Dayan:91}, enabling the highest possible separability of binary vectors during retrieval; this, in effect, maximizes the signal-to-noise ratio. Greater levels of separability can be achieved when restricting retrieval to solely pairwise dependencies, thereby incurring quadratic computational complexity as observed in attention mechanisms~\citep{Krotov:16,Krotov:17,Ramsauer:21}. Notably, the covariance update is tantamount to the mechanism utilized in Fast Weight Programmers \citep{Schmidhuber:92ncfastweights,Schlag:21}, which were subsequently augmented with a fixed decay factor applied to $\mathbf{C}_{t-1}$ as well as a static learning rate that scales 
$\mathbf{v}_{t} \mathbf{k}_t ^\top$~\citep{Ba:16}.
Following this paradigm, we embed the covariance update process within the LSTM structure, aligning the forget gate with the decay factor, associating the input gate with the learning rate, and employing the output gate to modulate the resulting retrieved vector.

For this matrix memory, the normalizer state is the weighted sum of key vectors, where each key vector is weighted by the input gate and all future forget gates. Again, the normalizer state keeps record of the strength of the gates. Since the dot product between query and normalizer state can be close to zero, we use the absolute value of this dot product and lower bound it by a threshold (typically 1.0) as done previously~\citep{Sun:23arxiv}. The mLSTM forward pass is:

\begin{align}
\label{eq:mlstm_recurrent_begin}
\mathbf{C}_t \ &= \  f_t \ \mathbf{C}_{t-1} \ + \ 
  i_t \ \mathbf{v}_t \ \mathbf{k}_t^\top & &  &\text{cell state} \\
\mathbf{n}_t \ &= \  f_t \ \mathbf{n}_{t-1} \ + \ 
  i_t \ \mathbf{k}_t & &  &\text{normalizer state} \\
\mathbf{h}_t  \ &= \ \mathbf{o}_t \ \odot \ \tilde{\mathbf{h}}_t
\ , \qquad \qquad 
\tilde{\mathbf{h}}_t \ = \   \mathbf{C}_t \mathbf{q}_t \ / \ 
  \max \left\{ |\mathbf{n}_t^\top \mathbf{q}_t|, 1 \right\} 
   &&&\text{hidden state} \\
\mathbf{q}_t \ &= \ \mathbf{W}_q \ \mathbf{x}_t \ + \ \mathbf{b}_q  & &  &\text{query input} \\
\mathbf{k}_t \ &= \ \frac{1}{\sqrt{d}} \mathbf{W}_k \ \mathbf{x}_t \ + \ \mathbf{b}_k  & &  &\text{key input} \\
\mathbf{v}_t \ &= \ \mathbf{W}_v \ \mathbf{x}_t \ + \ \mathbf{b}_v  & &  &\text{value input} \\
i_t \ &= \ \!\exp \!  \! \left( \tilde{i}_t  \right) \ , 
 \qquad \qquad \ \ \ \ \,
  \tilde{i}_t \ = \ \mathbf{w}^\top_{i} \ \mathbf{x}_t \ + \  b_{i} \ \ \ 
  &&&\text{input gate} \\
f_t \ &= \ \sigma \!  \left(  \tilde{f}_t \right) \ \text{OR} \ \!\exp \!  \! \left(  \tilde{f}_t \right) , \ \ \: \!
\tilde{f}_t \ = \ \mathbf{w}^\top_{f} \ \mathbf{x}_t  \ + \
  b_{f}
  &&& \text{forget gate} \\
\label{eq:mlstm_recurrent_end}
\mathbf{o}_t \ &= \ \sigma \left( \tilde{\mathbf{o}}_t \right) \ , \qquad \qquad \quad \ \ \ \,
  \tilde{\mathbf{o}}_t  \ = \ \mathbf{W}_{\mathbf{o}} \ \mathbf{x}_t \ + \
  \mathbf{b}_{\mathbf{o}}
  &&&\text{output gate} 
\end{align}

mLSTM is capable of incorporating several memory cells similar to the traditional LSTM architecture. In mLSTM, having numerous heads or numerous cells is functionally the same because memory mixing does not occur. To ensure the exponential gates in mLSTM remain stable, we employ stabilization methods identical to those used in sLSTM (refer to Equation~\ref{eq:slstmstabil}).  
Due to the absence of memory mixing in mLSTM, its recurrence relation lends itself to a parallelized reformulation. Further information is available in Appendix~\ref{sec:appmLSTM}.

\subsection{xLSTM Architecture}
\label{sec:xLSTMarch}

\paragraph{xLSTM Blocks.}
\label{sec:xLSTMblock}
The primary function of an xLSTM block is to produce a non-linear representation of previous sequence elements in a high-dimensional space, enabling the model to more effectively distinguish among various historical sequences or contexts. Successfully differentiating among distinct histories is essential for accurate prediction of subsequent sequence elements, such as the next token in a sequence. This approach leverages Cover’s Theorem~\citep{Cover:65}, which asserts that non-linearly transformed patterns in higher-dimensional spaces are more amenable to linear separation than those in their initial space. We investigate two residual block designs: (i) A residual block employing a post up-projection—similar to Transformer architectures—first condenses the history non-linearly in the original space, then uses a linear transformation to project into a higher-dimensional space, applies a non-linear activation, and finally projects linearly back to the original space; refer to the left panel of Figure~\ref{fig:Backbones} and the third column of Figure~\ref{fig:xlstm_sketch}. Further architectural details are provided in Figure~\ref{fig:appsLSTM_detailed} in the appendix. (ii) A residual block with pre up-projection—reminiscent of State Space Models—begins by linearly projecting into a high-dimensional space,

The model encodes historical information using a nonlinear function within a high-dimensional latent space, followed by a linear transformation that projects the data back into its original dimensionality. When an xLSTM block incorporates an sLSTM, a post up-projection mechanism is applied. Conversely, if the block integrates an mLSTM, the pre up-projection is employed due to the increased memory resources available within the high-dimensional space. Further clarification can be found in the left panel of Figure~\ref{fig:Backbones}, the third column of Figure~\ref{fig:xlstm_sketch}, and Figure~\ref{fig:appsLSTM_detailed} located in the appendix.

\paragraph{xLSTM Architecture. }
\label{sec:xLSTMarchitecure}
The xLSTM framework is formulated by sequentially stacking component modules in a residual manner, as outlined in prior works~\citep{Srivastava:15,He:16}. The backbone follows the prevalent pre-LayerNorm~\citep{Ba:16b} scheme, consistent with practices adopted in modern Large Language Models. Illustrative details are available in the final column of Figure~\ref{fig:xlstm_sketch}.

\subsection{Memory and Speed Considerations}

In contrast to Transformers, xLSTM networks exhibit linear computational requirements and their memory usage remains constant regardless of sequence length. Due to its compressive memory structure, xLSTM is particularly advantageous for deployment in industrial settings and for applications at the edge.

The memory used by mLSTM operates without learnable parameters; however, it demands substantial computation because of its $d \times d$ matrix memory and corresponding update operations. This design intentionally balances between maximizing memory capacity and managing computational cost. Nevertheless, these computations are amenable to parallelization on GPUs, resulting in only negligible effects on actual processing time.

Although mLSTM can be parallelized—similar to FlashAttention~\citep{Dao:22,Dao:23} or GLA~\citep{Yang:23arxiv}—the sLSTM architecture is inherently sequential and thus not parallelizable due to its use of memory mixing via hidden-to-hidden interactions. Despite this limitation, we engineered an efficient CUDA-based implementation, incorporating GPU memory optimizations down to the register, ensuring that the execution time is usually less than double that of mLSTM.

\section{Related Work}
\label{sec:related}

\paragraph{Linear Attention.} A variety of strategies have been proposed to address the Transformer’s quadratic scaling with respect to context length, enabling attention to operate with linear complexity. The Synthesizer introduces attention weights synthetically, omitting direct interactions between tokens~\citep{Tay:20arxiv}. Linformer employs a low-rank matrix to construct self-attention and achieves a linear approximation of the operation~\citep{Wang:20arxiv}. Linear Transformer redesigns the attention formulation to allow for linear computation~\citep{Katharopoulos:20}. Performer leverages positive orthogonal random features to linearly estimate the softmax attention~\citep{Choromanski:21}. Furthermore, methods such as Structured Global Convolution (SGConv)~\citep{Li:22} and Hyena Hierarchy~\citep{Poli:23} have substituted attention with efficient long-range convolutional operations.

\paragraph{State Space Models.} In recent times, State Space Models (SSMs) have gained considerable attention owing to their linear scalability with respect to context length and their competitive results against Transformer architectures. Early notable developments in this direction include the Structured State Space sequence model (S4)~\citep{Gu:21}. Subsequent advancements introduced several variants such as the Diagonal State Space (DSS) model~\citep{Gupta:22}, Gated State Space (GSS) models~\citep{Mehta:22}, S5 model~\citep{Smith:22}, Bidirectional Gated SSM (BiGS)~\citep{Wang:22}, H3 model~\citep{Fu:23}, as well as Mamba~\citep{Gu:24arxiv}.

\paragraph{Recurrent Neural Networks.} Recurrent Neural Networks (RNNs) have been proposed as alternatives to Transformer models and attention mechanisms, owing to their computational efficiency that scales linearly with sequence length. Advances such as Deep Linear Recurrent Units (LRUs) have demonstrated notable effectiveness in language modeling tasks~\citep{Orvieto:23, De:24arxiv}. Similarly, Hierarchically Gated Linear RNN (HGRN)~\citep{Qin:23} and its successor HGRN2~\citep{Qin:24arxiv} have also delivered strong outcomes. Among approaches for large language models, the RWKV method~\citep{Peng:23arxivshort,Peng:24arxivshort} has gained attention for achieving performance levels comparable to that of Transformers.

\paragraph{Gating.}
Gating, a central concept originating from LSTM, has subsequently been revisited and adapted in a variety of modern architectures. This mechanism plays a role in HGRN~\citep{Qin:23}, HGRN2~\citep{Qin:24arxiv}, Gated Linear Attention (GLA)~\citep{Yang:23arxiv}, Gated State Space (GSS) models~\citep{Mehta:22}, Bidirectional Gated SSM (BiGS)~\citep{Wang:22}, Moving Average Equipped Gated Attention (MEGA)~\citep{Ma:22}, RWKV~\citep{Peng:23arxivshort}, and Mamba~\citep{Gu:24arxiv}.

\paragraph{Covariance Update Rule.} To improve the memory capacity of the mLSTM cell, we augmented it with a matrix memory utilizing a covariance update rule. Comparable strategies are seen in Fast Weight Programmers \citep{Schmidhuber:92ncfastweights,Schlag:21}, RWKV-5 and RWKV-6~\citep{Peng:24arxivshort}, Retention~\citep{Sun:23arxiv}, Linear Transformer~\citep{Katharopoulos:20}, as well as HGRN2~\citep{Qin:24arxiv}, which all incorporate similar update procedures.

\paragraph{Most Related.} The models that are most conceptually akin to xLSTM include Retention~\citep{Sun:23arxiv}, RWKV~\citep{Peng:23arxivshort,Peng:24arxivshort}, and HGRN2~\citep{Qin:24arxiv}. These models employ mechanisms such as matrix memory and/or gating. In comparison to the recently introduced sLSTM, these models lack the capacity for memory mixing. This capability is crucial for resolving state tracking challenges, thereby granting LSTMs greater expressive power than State Space Models (SSMs) and Transformers~\citep{Merrill:24,Deletang:23}. Accurate state tracking is indispensable for tasks like code execution or maintaining entity information throughout extended narratives.

\paragraph{Residually Stacking Architectures.} The xLSTM models, in keeping with most recent advances in deep learning, employ a design in which their fundamental units are stacked with residual connections~\citep{Srivastava:15,He:16}.

This methodology has been instrumental in facilitating the development of very deep convolutional networks~\citep{He:16} and also forms the basis for Transformer architectures~\citep{Vaswani:17}. Transformers have become foundational to Large Language Models (LLMs) such as GPT-3~\citep{Brown:20short}, ChatGPT~\citep{Schulman:22short}, GPT-4~\citep{Achiam:23gpt4short}, Megatron-LM~\citep{Shoeybi:19}, Gopher~\citep{Rae:21short}, ERNIE 3.0 Titan~\citep{Wang:21arxivshort}, GLaM~\citep{Du:21glamshort}, the Chinese M6~\citep{Lin:21}, the multilingual AlexaTM 20B~\citep{Soltan:22}, OPT~\citep{Zhang:22opt}, Chinchilla~\citep{Hoffmann:22short}, BLOOM~\citep{Scao:22arxivshort}, GLM-130B~\citep{Zeng:22arxivshort}, LaMDA~\citep{Thoppilan:22short}, PaLM~\citep{Chowdhery:22short}, Llama~\citep{Touvron:23arxiv}, and Gemini~\citep{GeminiTeam:23arxiv,Reid:24arxiv}.

\section{Experiments}
\label{sec:Exp}

This section presents a systematic empirical analysis of xLSTM in comparison to prevailing approaches, with an emphasis on language modeling performance. Section~\ref{sec:ExpSyntethic} examines xLSTM's functionality through controlled synthetic tasks. In Section~\ref{sec:ExpComparison}, we benchmark the validation perplexity achieved by several state-of-the-art language modeling models, all trained on a dataset comprising 15B tokens from SlimPajama~\citep{Soboleva:23}. Additionally, we conduct ablation studies targeting various components of xLSTM on the same training corpus. The scaling properties of each model are then analyzed, following methodologies similar to those of \citet{Kaplan:20} and \citet{Brown:20short}.

Section~\ref{sec:ExpLanguage} advances to a comprehensive language modeling investigation. Here, both xLSTM and the leading models identified in Section~\ref{sec:ExpComparison} are trained on 300B SlimPajama tokens~\citep{Soboleva:23}. First, we evaluate the capacity of the models to generalize to much longer input sequences. Second, performance is measured by both validation perplexity and effectiveness on a range of downstream tasks~\citep{Sutawika:24}. Third, evaluation extends to 571 distinct text domains within the PALOMA language benchmark suite~\citep{Magnusson:23arxivshort}. Finally, we revisit scaling analysis for all methods, this time using a dataset twenty times larger than the initial comparison.

\label{sec:xlstm_blocks_notation}
Throughout all experimental settings, we denote xLSTM[$a$:$b$] as specifying the proportion $a/b$ of xLSTM blocks constructed using mLSTM compared to those using sLSTM. For instance, xLSTM[7:1] refers to an architecture comprising eight total blocks, of which seven employ mLSTM and one utilizes sLSTM. If the total number of blocks is set to 48, this configuration yields 6 blocks based on sLSTM and 42 built upon mLSTM. Additionally, all of our experiments employ up-projection blocks positioned before and after for mLSTM and sLSTM, respectively.

\subsection{Synthetic Tasks and Long Range Arena}
\label{sec:ExpSyntethic}
To begin, we evaluate the performance of xLSTM's novel exponential gating combined with memory mixing in the context of formal language tasks~\citep{Deletang:23}. Next, we investigate how well xLSTM's newly introduced matrix memory operates in the Multi-Query Associative Recall scenario~\citep{Arora:23arxiv}. Lastly, we analyze xLSTM's capability to handle extended input sequences by measuring its results on the Long Range Arena benchmark~\citep{Tay:21}.

\paragraph{Evaluation of xLSTM's Exponential Gating with Memory Mixing.} We evaluate the performance of xLSTM equipped with exponential gating and memory mixing mechanisms, which are expected to facilitate solving state-tracking tasks~\citep{Merrill:24, Merrill:23}. To this end, we expand upon the formal language benchmarks introduced by \citet{Deletang:23}, adapting them for multi-length training to enable assessment of the model's extrapolation ability across sequence lengths. Comprehensive descriptions of each task and additional results are provided in Appendix~\ref{sec:appExpSynthetic-formalLang}. In our comparison, we assess xLSTM alongside various architectures such as Transformers, State Space Models, and standard Recurrent Neural Networks. Accuracy is reported based on target tokens relevant to the corresponding task, with scores normalized between 0 (chance performance) and 1 (perfect accuracy). We consider 2-block variants for the following methods: xLSTM[0:1] (sLSTM only), xLSTM[1:0] (mLSTM only), xLSTM[1:1], as well as Llama, Mamba, RWKV, Retention, Hyena, LSTM, and LSTM integrated into Transformer blocks (LSTM (Block)). The comparative results are illustrated in Figure~\ref{fig:formal-main}. It should be noted that architectures such as Transformers and State Space Models lacking memory mixing capabilities (and thus, unable to perform explicit state tracking) fail to solve, for example, regular grammar tasks such as parity detection. This outcome corroborates prior findings indicating that Transformers and State Space models are inherently less expressive than RNNs for such problems~\citep{Merrill:24, Merrill:23, Deletang:23}.

\paragraph{Test of xLSTM's Memory Capacities on Associative Recall Tasks.} In this study, we investigate the matrix-based memory mechanisms of xLSTM through the Multi-Query Associative Recall task~\citep{Arora:23arxiv}, focusing on measuring memory capacity. For each input sequence, a set of key-value pairs is randomly sampled from an extensive vocabulary, requiring the model to store and later retrieve these associations. To increase the challenge beyond the standard setup, we expand both the number of key-value pairs, reaching up to 256, and the context length, pushing it to 2048. This extended configuration provides a more comprehensive evaluation of the models' memory abilities. We benchmark the 2-block variants of Llama, Mamba, RWKV-5, RWKV-6, xLSTM[1:1], and xLSTM[1:0], assessing them based on their retrieval accuracy. Transformers like Llama—owing to their exponentially scaled memory in the coding dimension~\citep{Ramsauer:21}—serve as the reference for optimal performance. Detailed outcomes are provided in Figure~\ref{fig:mqar-main}. Among the models that are not Transformer-based, xLSTM[1:1] demonstrates superior performance, including when the model size is small. Notably, the inclusion of the sLSTM block does not reduce memory capacity; on the contrary, it appears to enhance it, especially noticeable when handling the most challenging scenario involving 256 key-value associations. Further findings are given in Appendix~\ref{sec:appExpSynthetic-mqar}, where extrapolation analyses demonstrate that the improved memory of xLSTM enables it to generalize to sequence lengths surpassing those encountered during training.

\paragraph{Test of xLSTM's Long Context Capabilities on Long Range Arena.} To evaluate how xLSTM manages extended sequences and substantial context sizes, we conducted comparisons with alternative approaches using the Long Range Arena~\citep{Tay:21}. Across all tasks, xLSTM reliably achieves high performance, indicating that its architectural design effectively addresses various challenges presented by long context scenarios.

For more details, see Appendix~\ref{sec:appExpSynthetic-lra}.

\subsection{Method Comparison and Ablation Study}
\label{sec:ExpComparison}

To investigate our central inquiry—specifically, the capabilities of the newly proposed LSTM variants when applied to large-scale language modeling—we conduct training experiments involving xLSTMs, Transformers, State Space Models, and additional approaches using 15B tokens from SlimPajama under a consistent auto-regressive paradigm. Subsequently, we evaluate these models on a validation dataset and carry out ablation analyses focused on the xLSTM architectures.

\paragraph{Comparison Between xLSTM and Alternative Approaches.} To facilitate a robust comparison, all models are trained on a dataset comprising 15B tokens drawn from SlimPajama~\citep{Soboleva:23}. Their performance is measured by evaluating perplexity on a dedicated validation set. The suite of methods considered encompasses: xLSTM (our proposed approach), GPT-3 (Transformer)~\citep{Brown:20short}, Llama (Transformer)~\citep{Touvron:23arxiv}, H3 (SSM)~\citep{Fu:23}, Mamba (SSM)~\citep{Gu:24arxiv}, RWKV-4 (RNN)~\citep{Peng:23arxivshort}, RWKV-5 (RNN)~\citep{Peng:24arxivshort}, RWKV-6 (RNN)~\citep{Peng:24arxivshort}, GLA (linear Transformer)~\citep{Yang:23arxiv}, HGRN (RNN)~\citep{Qin:23}, HGRN2 (RNN)~\citep{Qin:24arxiv}, RetNet (linear Transformer)~\citep{Sun:23arxiv}, Hyena (linear Transformer)~\citep{Poli:23}, xLSTM[1:0], and xLSTM[7:1] (refer to Section~\ref{sec:xlstm_blocks_notation} for notation). 
All models employ mixed precision during training, except RWKV-5, RWKV-6, GLA, and HGRN2, for which the process does not use PyTorch's built-in automated mixed precision (see also Appendix Section~\ref{sec:appGenTrainProc}).
We organize the methods according to their architectural type: (a) Transformer models, (b) State Space Models (SSMs), and (c) Recurrent Neural Networks (RNNs) as well as linear Transformers. Linear Transformers adopt a linear alternative to the standard Transformer attention. Each architecture has a parameter count comparable to GPT-3 at 350M, with an embedding dimension of 1024 and 24 residual blocks. It is worth noting that only GPT-3 applies shared weights to both token and output embeddings, yielding a lower overall parameter count. The results presented in Table~\ref{tab:spaj15b_model_results} reveal that xLSTM achieves the best validation perplexity across all approaches. Further details are provided in Appendix~\ref{sec:appExpComparison}.
The scaling behavior observed in Figure~\ref{fig:temp_sclaw15B} suggests xLSTM maintains its advantage as model size increases.

\paragraph{Ablation Studies.}
Table~\ref{tab:spaj15b_model_results} and Figure~\ref{fig:temp_sclaw15B} reveal that xLSTM delivers strong language modeling performance when trained with 15B tokens drawn from SlimPajama. To systematically isolate the impact of architectural modifications, we begin with a standard LSTM model and incrementally convert it to xLSTM. Our first adjustment introduces LSTM layers into pre-LayerNorm residual structures. Next, we adapt the setup by incorporating a block following up-projection. In the final phase, exponential gating and matrix memory mechanisms are added.

The results are shown in Table~\ref{tab:ablstudies} (top). 

The results from the ablation studies suggest that the notable gains in performance stem from the combined effects of exponential gating and matrix memory. Recognizing the critical role of gating in RNNs and State Space Models, various gating strategies are systematically evaluated. As shown in Table~\ref{tab:ablstudies} (bottom), the findings demonstrate that making each gate both trainable and responsive to the input leads to progressively better outcomes. Further detailed examinations of each backbone component can be found in Appendix~\ref{sec:appAblDetails}.

\subsection{xLSTM as Large Language Model}
\label{sec:ExpLanguage}

The final phase of this research involves comprehensive language modeling experiments aimed at evaluating the viability of xLSTM as a large language model (LLM). Accordingly, we scale up the dataset, training xLSTM on 300B tokens sourced from SlimPajama. This token count aligns with that utilized in prior works such as Mamba~\citep{Gu:24arxiv} and Griffin~\citep{De:24arxiv}.

Our comparative analysis involves xLSTM, RWKV-4, Llama, and Mamba, selecting one representative from each category of methods as defined in Section~\ref{sec:ExpComparison}.
RWKV-4 is chosen to represent the RNN category because acceptable training precision configurations for RWKV-5, RWKV-6, and HGRN2 were only determined after initiation of the 300B token experiments (refer to Appendix~\ref{sec:appGenTrainProc}).
We employ multiple model scales—125M, 350M, 760M, and 1.3B parameters—and systematically assess length extrapolation, validation performance, and downstream task outcomes. Additionally, we analyze their efficacy in language modeling across 571 text domains included in the PALOMA benchmark, and conclude by exploring their scaling law properties.

\paragraph{Sequence Length Extrapolation.}
\label{sec:spaj300B_ctx_extrapolation}
We begin by evaluating the sequence length extrapolation capabilities of the 1.3B-parameter variants of xLSTM, RWKV-4, Llama, and Mamba. Although training is conducted using a context length of 2048, inference experiments are carried out with context sizes scaling up to 16384. Refer to Figure~\ref{fig:spaj300B_seq_extrapolation} for a summary of these findings. Notably, xLSTM models retain lower perplexity across expanded contexts compared to alternative approaches.

\paragraph{Validation Perplexity and Downstream Tasks.}
\label{sec:spaj_downstream_eval}
Next, across every model scale, we measure xLSTM, RWKV-4, Llama, and Mamba models on next token prediction with the SlimPajama validation set, as well as on downstream benchmarks targeting common sense reasoning capabilities.

The third column in Table~\ref{tab:lmeval} displays validation set perplexity scores for various methods. Across all model scales, xLSTM[1:0] and xLSTM[7:1] consistently outperform other models in terms of validation set perplexity. The additional columns in Table~\ref{tab:lmeval} summarize results from downstream tasks. With few exceptions, xLSTM ranks highest for nearly all tasks and model sizes; the ARC task is the only case where Mamba occasionally achieves superior performance. Further specifics can be found in Appendix~\ref{sec:appExpLanguage}.

\paragraph{Performance on PALOMA Language Tasks.} Next, we evaluate the next token prediction accuracy across all model sizes for xLSTM, RWKV-4, Llama, and Mamba on the PALOMA language tasks~\citep{Magnusson:23arxivshort}. Our metric for comparison is perplexity in next token prediction, calculated over 571 distinct text domains, which include sources as varied as nytimes.com and Reddit’s r/depression.

As summarized in Table~\ref{tab:ppleval}, token prediction perplexity is organized into broader language modeling categories (first seven columns) and more detailed domain-specific benchmarks (final five columns). On these tasks, xLSTM[1:0] demonstrates superior performance to xLSTM[7:1].

Specifically, xLSTM[1:0] achieves lower perplexity than Mamba in 568 out of 571 text domains (99.5\%), outperforms Llama in 486 domains (85.1\%), and surpasses RWKV-4 in 570 domains (99.8\%). Further results can be found in Appendix~\ref{sec:appExpLanguage}.

\paragraph{Scaling Laws.} Next, we analyze the power-law scaling characteristics, which provide predictive insights into performance as model size increases~\citep{Kaplan:20,Brown:20short}. Figure~\ref{fig:temp_sclaw300B} illustrates this scaling trend. Although all models exhibit analogous scaling patterns, their performance levels differ. RWKV-4 shows the lowest performance, with Llama and Mamba following sequentially. Notably, xLSTM achieves superior results compared to Mamba, maintaining a comparable advantage as Mamba holds over Llama. This scaling analysis suggests that as models become larger, xLSTM will likely maintain its advantageous performance in relation to both Transformers and State-Space models.

\textbf{Generation Times and Maximal Throughput.} We report the measured text generation latency in Figure~\ref{fig:inference_time_generation}, as well as the peak throughput results in Figure~\ref{fig:inference_time_decoding_throughput} (left), focusing on xLSTM configurations at the 1.3B parameter setting. Our evaluation includes similarly sized implementations of Mamba, Llama, and RWKV sourced from HuggingFace, with the Llama model utilizing a static key-value cache. At the time we performed these experiments, it was not possible to successfully compile the complete cache for the Transformer Model, nor to use \texttt{torch.compile} with the Mamba model. Each text generation benchmark is conducted at batch size 1 and with a pre-fill value of 16, which notably favors the Transformer model. 

As illustrated in Figure~\ref{fig:inference_time_generation}, xLSTM along with other recurrent models such as Mamba and RWKV-4 demonstrate linear complexity with sequence length, contrasting with the quadratic complexity observed for Llama. When assessing decoding throughput, we evaluate a spectrum of batch sizes and prefill values specifically for the Llama implementation. The analysis presented in Figure~\ref{fig:inference_time_decoding_throughput} (right) indicates that, due to its fixed memory requirements, xLSTM supports significantly larger batch sizes than Llama, leading to superior throughput performance.

\section{Limitations}

(i) Unlike mLSTM, the sLSTM’s approach to memory mixing inhibits the use of parallel operations, preventing efficient parallel processing. Despite this constraint, we created a dedicated CUDA kernel for sLSTM that currently executes at less than twice the latency of our parallelized mLSTM kernel. (ii) The CUDA kernels employed for mLSTM lack advanced optimizations; as a result, the present mLSTM implementation runs at roughly a quarter of the speed compared to FlashAttention or the scan method incorporated in Mamba. Performance could be improved significantly by designing faster CUDA kernels following techniques similar to those in FlashAttention. (iii) Processing the matrix memory in mLSTM incurs substantial computational expense, as operations must handle $d \times d$ matrices. However, since memory update and retrieval are parameter-free and leverage conventional matrix-based parallelization, the real-world computational delay from this memory complexity remains relatively negligible. (iv) Careful selection of initialization values for the forget gates is essential. (v) The matrix memory does not depend on the sequence length, meaning that increasing the sequence size can strain the memory when contexts get longer. Nevertheless, for contexts of up to 16k tokens, this effect does not constitute a practical bottleneck, as shown in Section~\ref{sec:spaj300B_ctx_extrapolation}.
(vi) Owing to the significant computational requirements of large-scale language modeling experiments, we refrained from both full architectural optimization and thorough hyperparameter tuning, particularly with the larger xLSTM models. We expect that more rigorous optimization will be necessary for xLSTM to realize its full capabilities.

\section{Conclusion}
We have addressed, in part, the primary question: How effective is large-scale LSTM for language modeling when extended to billions of parameters? Our findings indicate that such models are at minimum comparable to leading methods including Transformers and State Space Models. The introduction of xLSTM—incorporating exponential gating through memory mixing and a novel memory architecture—represents a significant enhancement. When evaluated against the latest approaches like Transformers and State Space Models, xLSTM demonstrates superior or equivalent performance in language modeling tasks. The observed scaling behavior suggests that expanding xLSTM could position it as a formidable challenger to current Large Language Models employing Transformer architectures. Moreover, xLSTM holds promise for substantial influence in diverse areas such as Reinforcement Learning, Time Series Prediction, and the modeling of physical phenomena.

\end{document}